\section{Solusi Schwarzschild}
\label{sec:solusi-schwarzschild}

\noindent Persamaan medan Einstein pada Persamaan~\eqref{eq:efe} berlaku sangat umum untuk 
sembarang distribusi materi dan energi, sehingga secara matematis sulit diselesaikan secara 
langsung. Kasus paling sederhana yang dapat ditinjau adalah kondisi vakum ($T_{\mu\nu} = 0$), 
yakni ketika ditinjau geometri ruang-waktu di luar suatu objek masif di mana tidak terdapat 
materi maupun energi. Pada kondisi ini, ruas kanan persamaan medan Einstein lenyap. Dengan mengambil trace 
persamaan medan Einstein pada kondisi vakum, diperoleh bahwa skalar Ricci $R$ juga 
lenyap, sehingga persamaan medan Einstein pada kondisi vakum tereduksi menjadi 
$R_{\mu\nu} = 0$. Perlu ditekankan bahwa lenyapnya tensor Ricci 
tidak berarti ruang-waktu menjadi datar, karena tensor Riemann yang memuat seluruh informasi 
kelengkungan masih dapat bernilai tak nol; ruang-waktu vakum di sekitar objek masif tetap 
melengkung meskipun tidak terdapat materi di titik yang ditinjau \parencite{Carroll2004}. 
Penyelesaian persamaan medan Einstein pada kondisi vakum inilah yang pertama kali diperoleh 
sebagai solusi eksak oleh Karl Schwarzschild pada tahun 1916, tidak lama setelah Einstein 
mempublikasikan teori relativitas umum \parencite{Schwarzschild1916}.

Untuk menyelesaikan persamaan medan Einstein pada kondisi vakum, Schwarzschild mengasumsikan bahwa ruang-waktu di luar objek masif yang tidak berotasi dan tidak bermuatan bersifat statis dan bersimetri bola \parencite{Wald1984,Carroll2004}. Sifat statis menunjukkan bahwa geometri ruang-waktu tidak berubah terhadap waktu, sedangkan simetri bola menunjukkan bahwa geometri tetap sama terhadap rotasi di sekitar pusat objek.

Dengan menggunakan koordinat bola $(t,r,\theta,\varphi)$, kedua asumsi tersebut memungkinkan elemen garis dituliskan dengan dua fungsi radial yang belum diketahui, masing-masing untuk komponen waktu dan komponen radial. Penyelesaian persamaan medan Einstein pada kondisi vakum untuk ansatz tersebut menunjukkan bahwa kedua fungsi ini saling berkebalikan \parencite{Wald1984}, sehingga elemen garis dapat dituliskan dengan satu fungsi metrik tunggal $f(r)$, menghasilkan solusi Schwarzschild
\begin{equation}
ds^2=-f(r)\,dt^2+f(r)^{-1}\,dr^2+r^2\left(d\theta^2+\sin^2\theta\,d\varphi^2\right),
\label{eq:metrik-schwarzschild}
\end{equation}
dengan $f(r)$ merupakan fungsi metrik yang bergantung pada koordinat radial. Bagian angular $r^2(d\theta^2+\sin^2\theta\,d\varphi^2)$ menunjukkan sifat simetri bola, sedangkan $r$ menyatakan radius. Bentuk fungsi metrik yang memenuhi persamaan medan Einstein pada kondisi vakum diberikan oleh
\begin{equation}
f(r)=1-\frac{2GM}{c^2r}.
\label{eq:fungsi-metrik-si}
\end{equation}

Dalam \textit{natural unit} dengan $G=c=1$, fungsi metrik tersebut menjadi
\begin{equation}
f(r)=1-\frac{2M}{r}.
\label{eq:fungsi-metrik-schwarzschild}
\end{equation}
sehingga kesamaan dimensi antara $r$ (panjang) dan $M$ (massa) yang tampak pada persamaan 
ini semata-mata merupakan konsekuensi dari pemilihan satuan tersebut, bukan berarti kedua 
besaran tersebut secara fisis setara. Fungsi metrik inilah yang sepenuhnya menentukan 
struktur geometri ruang-waktu Schwarzschild \parencite{Wald1984}.

Solusi Schwarzschild mendeskripsikan geometri ruang-waktu di luar sembarang 
objek bermassa $M$ yang tidak berotasi dan tidak bermuatan, selama daerah yang ditinjau 
berada di luar permukaan objek tersebut. Pada objek astrofisika seperti bintang atau planet, 
daerah di dalam permukaan objek masih mengandung materi, sehingga kondisi vakum tidak 
terpenuhi dan solusi Schwarzschild tidak berlaku di sana; solusi ini hanya menggambarkan 
geometri pada daerah kosong di luar permukaannya. Semakin kecil radius suatu objek untuk 
massa $M$ yang tetap, semakin dekat pula daerah berlakunya solusi Schwarzschild terhadap 
pusat objek tersebut, sehingga cakupan geometri yang dapat dideskripsikan oleh solusi ini 
menjadi semakin luas.

Pada jarak yang jauh dari sumber gravitasi ($r \to \infty$), fungsi metrik 
pada Persamaan~\eqref{eq:fungsi-metrik-schwarzschild} mendekati nilai satu, sehingga 
ruang-waktu Schwarzschild bersifat datar secara asimtotik. Hal ini berarti geometri 
ruang-waktu mendekati ruang-waktu Minkowski pada jarak yang jauh dari objek masif, sehingga 
pengaruh gravitasi melemah dan menghilang seiring bertambahnya jarak. Sifat datar asimtotik 
inilah yang menjadi karakteristik utama ruang-waktu Schwarzschild dengan konstanta 
Kosmologi bernilai nol.

\subsection{Konstanta Kosmologi dalam Ruang-Waktu Schwarzschild}
\label{subsec:konstanta-Kosmologi}

\noindent Solusi Schwarzschild pada Persamaan~\eqref{eq:metrik-schwarzschild} 
diperoleh dengan mengasumsikan konstanta Kosmologi $\Lambda$ bernilai nol. Konstanta 
Kosmologi merupakan suku tambahan pada persamaan medan Einstein yang diperkenalkan oleh 
Einstein pada tahun 1917 dengan alasan meyakini alam semesta bersifat statis 
\parencite{Einstein1917}. Tanpa suku $\Lambda$, persamaan medan Einstein memprediksi bahwa 
alam semesta harus runtuh akibat tarikan gravitasi materi. Untuk mencegah hal ini, Einstein 
menambahkan $\Lambda$ sebagai suku penolak yang menciptakan keseimbangan, sehingga 
menghasilkan model alam semesta yang statis \parencite{Carroll2004}.

Ketika konstanta Kosmologi bernilai tak nol ($\Lambda \neq 0$), persamaan medan Einstein 
pada kondisi vakum tetap diselesaikan dengan asumsi ruang-waktu yang sama seperti sebelumnya, 
yaitu statis dan bersimetri bola, sehingga bentuk umum tensor metrik dan elemen garis tetap 
mengikuti Persamaan~\eqref{eq:metrik-schwarzschild}; hanya fungsi metriknya yang berubah. 
Penyelesaian persamaan menghasilkan fungsi metrik baru,
\begin{equation}
    f(r) = 1 - \frac{2M}{r} - \frac{\Lambda}{3} r^2,
    \label{eq:fungsi-metrik-umum}
\end{equation}
yang pertama kali diperoleh oleh Kottler \parencite{Kottler1918}, sehingga menggeneralisasi 
solusi Schwarzschild dengan menambahkan satu suku tambahan yang sebanding dengan $\Lambda r^2$. 
Kehadiran suku tambahan inilah yang mengubah perilaku asimtotik geometri ruang-waktu pada 
jarak yang jauh, sehingga nilai dan tanda dari $\Lambda$ akan sangat menentukan struktur 
geometri serta karakteristik ruang-waktu yang dihasilkan.

Dalam penelitian ini, fokus diberikan pada dua kasus khusus. Pada kasus $\Lambda = 0$, 
diperoleh ruang-waktu Schwarzschild yang datar secara asimtotik (Schwarzschild 
\textit{Asymptotically Flat}). Pada kasus $\Lambda < 0$, ruang-waktu bersifat anti-de Sitter 
secara asimtotik (Schwarzschild Anti-de Sitter) dengan kelengkungan negatif konstan pada 
jarak jauh.